\section{Polygonal area and quadrilaterals}
\label{sec:polygons}
A polygon is a finite ordered vertex list, closed by repeating its first
vertex. A convex polygon here has nonempty interior and lists exactly
the extreme vertices of its convex hull in counterclockwise boundary
order; redundant collinear vertices are omitted. A triangulated polygon is supplied
with a finite partition into nonoverlapping triangles: each internal edge
occurs twice in opposite directions and the remaining edges are its
oriented boundary. These conditions can be checked from exact coordinates
by orientation tests and a finite arrangement of segments. We do not
assume that an arbitrary self-intersecting vertex list bounds a simple region.

\begin{theorem}[Shoelace formula and finite dissection invariance]
\label{poly:shoelace}
If $P_0,\ldots,P_{n-1}$ is the boundary of a positively triangulated
polygon, its area is
\[
 \area(P)=\frac12\sum_{i<n}[P_i,P_{i+1}],\qquad P_n=P_0.
\]
The value is independent of the chosen positive triangulation, unchanged
by translating all vertices or subdividing boundary edges, and additive
under a finite partition into polygons with disjoint interiors.
\end{theorem}
\begin{proof}
The doubled oriented area of a triangle $ABC$ is
$[B-A,C-A]=[A,B]+[B,C]+[C,A]$. Sum this identity over the supplied
triangles. Every internal edge contributes $[X,Y]+[Y,X]=0$; all
remaining terms are exactly the boundary sum. Thus any two positive
triangulations have the same sum of unsigned triangle areas. Adding a
common translation $a$ changes each summand by
$[a,P_{i+1}]-[a,P_i]$, which telescopes to zero. Inserting
$Q=(1-t)P_i+tP_{i+1}$ replaces $[P_i,P_{i+1}]$ by
$[P_i,Q]+[Q,P_{i+1}]$, the same value. Summing the boundary identities
over a partition cancels all its internal edges too.

For a convex polygon the fan from $P_0$ is a positive triangulation.
Every ray from $P_0$ through the polygon meets its boundary in a segment;
the boundary vertices order those rays and the consecutive triangular
sectors have disjoint interiors. This is checked by their determinant
signs. The fan therefore gives $n-2$ triangles without any triangulation
search. For arbitrary supplied partitions the cancellation invariant is
the partial sum of processed directed edges, a quantifier-free equation
of coded lists, so the proof is uniform in the number of triangles.
\end{proof}

\begin{theorem}[Convex polygon angle sum]\label{poly:angles}
The interior angles of a convex $n$-gon sum to $(n-2)$ straight angles.
Its positive exterior turns sum to one full turn.
\end{theorem}
\begin{proof}
Use the convex fan in Theorem~\ref{poly:shoelace}. Each of its $n-2$
triangles contributes one straight angle by Theorem~\ref{tri:angle-sum}.
At each vertex the fan rays subdivide exactly that polygon's interior
angle, so adding the fan contributions gives the first assertion. Each
exterior turn is supplementary to its interior angle; subtracting from
$n$ straight angles leaves two straight angles.

This argument has a precise finite angle-code interpretation. Let
$H=\{(x,y):x^2+y^2=1,\ y>0\text{ or }(x,y)=(1,0)\}$.
Represent an angle by $(k,z)$, meaning $k$ half-turns and a remainder
$z\in H$. For $(k,z)+(l,w)$ form the pair product $p=z*w$.
If $p\in H$, return $(k+l,p)$; otherwise return $(k+l+1,-p)$.
The negative-axis endpoint increments the counter and the positive-axis
endpoint does not. Straight and full turns have canonical codes
$(1,(1,0))$ and $(2,(1,0))$. For remainders in $H$ the sum is
below two half-turns, so there is no missing double-carry case.
All tests are algebraic sign tests. The fan rays are ordered
by determinant signs and subdivide intervals of angle less than a
straight angle. Adding their adjacent sectors therefore has no unrecorded
wrap. Reassociating a finite sum changes neither its pair nor its carry:
refine the concatenated oriented sectors at both axis directions and
all participating endpoints. Both groupings count the same half-turn
crossings and retain the same final sector. The pair product is
associative by its coordinate identity, while the crossing count removes
the possible whole-turn ambiguity. This is a finite-list identity,
not equality merely modulo a full turn. The bounded sum invariant proves
the statement uniformly in $n$.
\end{proof}

\begin{theorem}[Areas of parallelograms, rectangles, triangles, and trapezoids]
\label{poly:areas}
A parallelogram on $u,v$ has area $|[u,v]|$; a triangle on the same
base vectors has half that area. A rectangle with perpendicular side
lengths $a,b$ has area $ab$. A trapezoid with parallel base lengths
$a,b$ and perpendicular height $h$ has area $(a+b)h/2$.
\end{theorem}
\begin{proof}
The parallelogram's diagonal divides it into the two equal determinant
triangles. If $u\cdot v=0$, the Gram identity gives
$[u,v]^2=N(u)N(v)$, so the nonnegative area is $ab$.
For a trapezoid choose an orthonormal frame with bases
$(0,0),(a,0)$ and $(t,h),(t+b,h)$, where $a,b,h>0$.
The boundary order is $(0,0),(a,0),(t+b,h),(t,h)$. Its shoelace
sum is $ah+bh$, independently of $t$. Orthogonal motions preserve
the absolute determinant, proving the general case. The same calculation
gives base times height divided by two for a triangle.
\end{proof}

\begin{theorem}[Equal-area figures between parallels]\label{poly:parallel-areas}
Triangles on the same base and between the same parallels have equal
area, as do parallelograms on the same base between the same parallels.
Conversely, two triangle vertices on the same side of a fixed nonzero
base give equal areas exactly when they lie on one parallel to that base.
\end{theorem}
\begin{proof}
For base $AB$ put $u=B-A$. Sliding a third vertex by $tu$ leaves
$[u,C-A]$ unchanged. On one fixed side the determinant sign is fixed,
so equality of unsigned areas is equality of these determinants.
Subtracting them gives $[u,D-C]=0$, precisely proportional directions.
Parallelogram area is twice the corresponding triangle area.
\end{proof}

\begin{theorem}[Parallelogram characterizations]\label{poly:parallelogram}
For a nondegenerate quadrilateral in boundary order, its diagonals bisect
each other if and only if $A+C=B+D$. This is equivalent to opposite
edge vectors being equal and parallel. In a parallelogram equal diagonals
characterize a rectangle, and perpendicular diagonals characterize a rhombus.
\end{theorem}
\begin{proof}
Equality of diagonal midpoints is precisely $A+C=B+D$. Rearrangement
gives $B-A=C-D$ and $D-A=C-B$, which are the two equal parallel
edge pairs. Translate $A$ to zero and write $B=u,D=v,C=u+v$.
The diagonal vectors are $u+v,v-u$.
Their squared-length difference is $4u\cdot v$, so equality means
adjacent sides perpendicular, the rectangle condition. Their dot product
is $N(v)-N(u)$, so perpendicularity means equal adjacent lengths,
the rhombus condition. Conversely each condition gives its diagonal
criterion by the same identities.
\end{proof}

\begin{theorem}[Varignon's midpoint parallelogram]\label{poly:varignon}
The successive side midpoints of any quadrilateral form a parallelogram,
possibly degenerate. For a convex quadrilateral its area is half the
original area. Its sides have half the lengths of the original diagonals.
\end{theorem}
\begin{proof}
Set $E=(A+B)/2$, $F=(B+C)/2$, $G=(C+D)/2$, $H=(D+A)/2$.
Then $F-E=(C-A)/2$, $G-F=(D-B)/2$, and the other two sides
are their negatives. The signed doubled area of the original cycle is
\[
 [A,B]+[B,C]+[C,D]+[D,A]=[C-A,D-B].
\]
The midpoint parallelogram area is
$|[(C-A)/2,(D-B)/2]|$, one quarter of this absolute determinant,
while the original convex area is one half of it. The vector identities
also give the side-length assertions.
\end{proof}

\begin{theorem}[Diagonal area formula and bound]\label{poly:diagonal-area}
For a convex quadrilateral with diagonal lengths $p,q$ and angle $\theta$
between their directions,
$\area=\tfrac12pq\sin\theta\le\tfrac12pq$.
Equality holds exactly when the diagonals are perpendicular.
\end{theorem}
\begin{proof}
The determinant identity in the preceding proof gives the area as
$|[C-A,D-B]|/2$. The definition of the angle pair identifies its
second coordinate with this determinant divided by $pq$.
The Gram identity bounds the determinant squared by $p^2q^2$,
with equality exactly when the diagonal dot product vanishes.
Both diagonals are nonzero, so the stated divisions are valid.
\end{proof}

\begin{theorem}[Bretschneider identity and the cyclic area maximum]
\label{poly:bretschneider}
Let a strictly convex quadrilateral have consecutive side lengths
$a,b,c,d$, area $K_0$, semiperimeter $s$, and opposite interior angles
$B,D$. Write their upper-semicircle pairs as $(x,u),(y,v)$.
Then
\[
 K_0^2=(s-a)(s-b)(s-c)(s-d)
          -\frac{abcd}{2}(1+xy-uv).
\]
The last term is nonnegative. Thus among convex quadrilaterals with these
side lengths, any cyclic realization has maximal area; equality in the
bound occurs exactly for a cyclic realization.
\end{theorem}
\begin{proof}
The diagonal $AC$ splits the quadrilateral into the triangles at $B,D$.
Put $X=ab$, $Y=cd$, and $E=a^2+b^2-c^2-d^2$.
The cosine law on that common diagonal gives $E=2Xx-2Yy$.
The area--sine formula gives $2K_0=Xu+Yv$. A direct difference of
squares factors as
\[
 16(s-a)(s-b)(s-c)(s-d)=4(X+Y)^2-E^2.
\]
Subtract $16K_0^2=4(Xu+Yv)^2$ and use $x^2+u^2=y^2+v^2=1$.
The result is $8XY(1+xy-uv)$, proving the identity after division by
16. Moreover
$2(1+xy-uv)=(x+y)^2+(u-v)^2\ge0$.
It vanishes exactly when $x=-y$ and $u=v$, meaning that the two
opposite angles are supplementary. The cyclic-angle theorem and its
converse from Section~\ref{sec:circles} say exactly that this happens
for a cyclic quadrilateral. All side lengths are positive, so multiplying
or cancelling $abcd$ does not lose a case.
\end{proof}

\begin{corollary}[Brahmagupta's formula]\label{poly:brahmagupta}
For a strictly convex cyclic quadrilateral,
\[
 \area=\sqrt{(s-a)(s-b)(s-c)(s-d)}.
\]
\end{corollary}
\begin{proof}
Its opposite angles are supplementary, so the correction term in
Theorem~\ref{poly:bretschneider} vanishes. Each factor is positive:
each side is shorter than the three-side broken path joining its endpoints,
by the strict triangle inequalities in the two triangles of a diagonal.
The area is positive, selecting the displayed nonnegative root uniquely.
\end{proof}

\begin{theorem}[Areas of similar figures on a right triangle]
\label{poly:similar-pythagoras}
Take any fixed positively triangulated polygonal shape. Similar copies
built with scale proportional to the three sides of a right triangle
have the area on the hypotenuse equal to the sum of the other two areas.
\end{theorem}
\begin{proof}
If the reference area is $A_0$ and the three scale factors are
$ka,kb,kc$, area scaling gives $A_0k^2a^2$, $A_0k^2b^2$,
$A_0k^2c^2$. Pythagoras supplies $c^2=a^2+b^2$.
The determinant proof of area scaling applies term by term to the finite
triangulation, so this includes any such shape, not only squares.
\end{proof}

\begin{theorem}[Regular polygon area and explicit elementary examples]
\label{poly:regular}
Let $n$ vertices on a circle of radius $r>0$ occur in convex boundary
order, with the same successive rotation pair $(c,t)$, $t>0$.
For $n\ge3$ its side length is $\ell=r\sqrt{2(1-c)}$, apothem is
$a=r\sqrt{(1+c)/2}$, and
\[
 \area=\frac{nr^2t}{2}=\frac{a(n\ell)}2.
\]
Regular triangles, squares, pentagons, and hexagons can be explicitly
constructed with straightedge and compass; repeated angle bisection
doubles any of these numbers of sides.
\end{theorem}
\begin{proof}
Adjacent radius vectors have dot product $r^2c$ and determinant $r^2t$.
The cosine law gives $\ell^2=2r^2(1-c)$. Their midpoint has squared
norm $r^2(1+c)/2$ and lies on the perpendicular to the chord; it is
the foot defining the apothem. Since $t>0$ and $c^2+t^2=1$,
$a\ell=r^2t$. Sum the $n$ oriented triangles from the center.

For a triangle use rotation pair $(-1/2,\sqrt3/2)$, for a square
$(0,1)$, and for a hexagon $(1/2,\sqrt3/2)$. Powers of these pairs
give respectively 3, 4, and 6 vertices and one full counterclockwise turn,
as their signs and products show. For a pentagon use
\[
 c=\frac{\sqrt5-1}{4},\qquad
 t=\frac{\sqrt{10+2\sqrt5}}4.
\]
Squaring verifies $c^2+t^2=1$. Put $f=4c^2+2c-1=0$ and $z=(c,t)$.
Expansion gives $\operatorname{Re}(z^5)-1=(c-1)f^2$ and
$\operatorname{Im}(z^5)=tf(4c^2-2c-1)$, so $z^5=(1,0)$.
Its successive nontrivial powers lie in quadrants
I, II, III, IV: $0<c<1/2$, the square has negative first and positive
second coordinate, and the last two powers are conjugates of the first
two. Hence their boundary order is the required five-step full turn,
not a star traversal. All displayed quantities use only arithmetic and
nonnegative square roots, whose straightedge--compass constructions were
proved in the circle chapter. The pair
$(\sqrt{(1+c)/2},\ t/(2\sqrt{(1+c)/2}))$ bisects a central angle
in $(0,\pi)$, so its successive powers insert one midpoint direction
between each original pair of radius directions. This doubles the side
count and proves the final construction claim.
\end{proof}
